\chapter{Science, physics and truth}
\label{SciencePhysPhilo}

\epigraph{No amount of experimentation can ever prove me right; a single experiment can prove me wrong.}{\textit{Albert Einstein}}

%Mention somewhere here that the objective of the text is to bridge the gap between the available popular literature and the actual study material which may be too demanding for people who do not have the ambition to study the topics thoroughly
\section{Distinguish fact from fiction}

\noindent As a young teenager, later than most of my peers, I realized that some statements about reality are true while others are not. Still other statements are simply undecidable (maybe because they are too vague). What I needed was a mechanism to be able to decide for myself which statement belonged to which category. It took me a couple of years more to realize that a method to do just that already existed: the scientific method.\\

\noindent So what is the scientific method? In its simplest form it amounts to making statements about the world, checking whether the statements hold in nature\footnote{By nature I mean all of reality, not just forests, mountains and the sea.} and getting rid of the ones that contradict what actually happens. A statement could be \itshape "all swans are white"\normalfont. We then travel the world in search of every swan we can find. Those we do find turn out all to be white. It seems "the experiment" corroborated the hypothesis. This does not mean the statement is 'true' as we may one day stumble upon a black swan we missed in our initial search, or a black swan may have been born since that time. All we can say until we encounter a black swan is that despite thorough trials to disprove the statement, it seems to hold true. Nevertheless, the search for a non-white swan continues. And if one day we indeed find such an unexpected black swan, we conclude that the most recent "experiment" falsified the hypothesis (which may even have become a respected theory by now). It is thrown in the garbage bin of historical ideas that may have been true but turned out not to be. Thus the scientific method filters out the statements which are untrue\footnote{The philosophy of science is much more refined and complex than is presented here and the given argument by no means does justice to the subtleties involved. The gist of it is nevertheless what most scientists and philosophers of science would call science.}.\\

\noindent So the objective is to describe reality, to find the patterns of nature, the laws of the world. Why then is there so much math involved? Would it not be possible to make statements like the one above without reverting to long mathematical formulas and calculations? After all, we did not need math to make a statement about swans nor to test its truth. There are in my opinion 3 main reasons for the ubiquity of math in science in general and physics in particular:
\begin{enumerate}
	
	\item \bfseries Precision of statements: \normalfont To assess whether a statement is true, it is important to understand what the statement is exactly. Is it clear what we mean by a swan? Do we know the exact characteristics of swans and how they differ from similar non-swan birds? Are we sure we know what we mean by white? Do they have to be all white? Would a white swan with a couple of freckles be sufficient to not call it white? The language in which a statement is presented needs to be clear and precise for it to be possible to determine whether it is true or false. Mathematics is just such a language that allows making very precise statements once the language is mastered.
	
	\item	\bfseries Quantification of statements: \normalfont As we look at all the swans we realize that they do not all have exactly the same color. They have different shades of white. We may then decide to make a whiteness scale. Every swan then becomes a mark on that scale. Somewhere on the scale we decide on the whiteness limit. Beyond that point, the swan cannot be considered white anymore. The statement has been quantified which allows for an even more precise analysis of the truth of the statement. 
	
	\item \bfseries Logical consequences of statements: \normalfont Math is more than merely a language. It is also a formalism which allows for deduction of the logical consequences of assumptions. Isaac Newton wrote down an equation which we will start discussing in chapter \ref{Intro_Physics}. That equation is meant to describe the motion of the moon, the flow of a river and the flight of an airplane. But going from his generic equation to the actual predictions for each case requires some work. Mathematics allows us to logically derive the particular from the general. Sometimes the derived consequences of assumptions even go against our intuition. For instance, if we assume that light travels at the same speed in vacuum, no matter who looks at it, it turns out that simultaneity breaks down. Two things that happen\footnote{"Things that happen" are often called "events" by physicists. If you want to fit in, you will need the jargon.} at the same time from one person's point of view will not necessarily happen at the same time from someone else's point of view. It is absolutely not trivial to see how an assumption about the speed of light can influence the flow of time or the order of events, but it is true nonetheless. Math allows us to think in a formal and logical way about the consequences of assumptions even when they do not seem intuitive or obvious.
\end{enumerate}
This text will focus mostly on physics. But it is unavoidable to introduce the math necessary to be able to understand physics beyond the popular literature. Of course, math has a beauty of its own, independently of its use in science. There will be short passages of math that are meant simply to show you some of her beauties and her surprises. But try to not get lost in the equations. They are there. They are important. But they are, despite the look of some physics textbooks, not the main point. They should be considered in this context merely as tools to convey our understanding of reality. This understanding is and always should be the focus when doing physics.\\

\noindent I will do my very best to make all of it as easy and approachable as it can be made. Nevertheless, you may (and probably will) find some parts challenging at a first read. That is completely normal. Just remember whenever you get stuck that learning is a gradual process and not having understood everything is not the same as having understood nothing at all. If at every reading of the text you gain a new insight, it will open doors to even further insights, it will allow you to ask new and sometimes deeper questions. And most importantly, it will hopefully give you that wonderful feeling of having understood something about the universe.